\documentclass[12pt]{article}
\usepackage{amsmath,amssymb,amsthm}
\usepackage{mathtools}
\usepackage{hyperref}
\usepackage{geometry}
\usepackage{booktabs}
\usepackage{array}
\geometry{margin=1in}
\hypersetup{colorlinks=true,linkcolor=blue,citecolor=blue}

\newtheorem{theorem}{Theorem}[section]
\newtheorem{lemma}[theorem]{Lemma}
\newtheorem{corollary}[theorem]{Corollary}
\newtheorem{proposition}[theorem]{Proposition}
\theoremstyle{definition}
\newtheorem{definition}[theorem]{Definition}
\newtheorem{remark}[theorem]{Remark}
\newtheorem{openproblem}[theorem]{Open Problem}
\newtheorem{observation}[theorem]{Computational Observation}

\newcommand{\ZZ}{\mathbb{Z}}
\newcommand{\ZK}{\mathbb{Z}/2^K\mathbb{Z}}
\newcommand{\ZKZ}{\mathbb{Z}_2}
\newcommand{\OK}{\Omega_K}
\newcommand{\vt}{v_2}
\newcommand{\TV}{\mathrm{TV}}
\newcommand{\Pk}{P_K}
\newcommand{\Pkm}{P_{K,m}}

\title{Scale-Coherence Rigidity and the Orbitwise Gap\\
in Collatz Dynamics\\[0.5em]
{\large Paper D: From Population Mixing to Individual Orbits}}

\author{Hiroki Kamanoi}
\date{Draft --- May 2026}

\begin{document}
\maketitle

\begin{abstract}
Papers B and C established exact population-level mixing for the Collatz dynamics modulo $2^K$: the distribution of $T^t(a) \bmod 2^K$ converges to uniform in $K-1$ steps (TV sense), and a single high-bit flip already creates a spectral gap $\delta_{K,1} \approx 0.29$ independent of $K$. The present paper addresses the fundamental gap between these \emph{population} statements and the \emph{individual orbit} statement required by the Collatz conjecture.

We identify three layers of structure. First, for any fixed orbit $\{n_t\}$, the effective high-bit parameter $m(t) = \lfloor\log_2 n_t\rfloor - K$ grows with $n_t$, making the K-level dynamics increasingly close to the true 2-adic Haar kernel---which, paradoxically, has \emph{exponentially small} spectral gap. Second, Weyl equidistribution implies that $n_t \bmod 2^K$ visits each residue with limiting frequency $1/n_K$, including residue $1$; but visiting $n_t \equiv 1 \pmod{2^K}$ does not imply $n_t = 1$. Third, the Collatz conjecture is equivalent to the \emph{simultaneous scale condition}: $n_t \equiv 1 \pmod{2^K}$ for \emph{all} $K$ simultaneously forces $n_t = 1$.

We formulate this as the \textbf{Simultaneous Scale Coherence Problem} and show that its resolution---which would imply the Collatz conjecture---requires exactly the multi-scale information rigidity identified in Papers B and C. We prove a conditional theorem: if the Open Gap Problem of Paper C holds ($\inf_K \delta_{K,1} > 0$), then the arithmetic memory requirements for a divergent orbit become increasingly contradictory across scales.
\end{abstract}

\tableofcontents\bigskip

%% ============================================================
\section{Introduction}
%% ============================================================

\subsection{The series so far}

This is the fourth paper in a series on the arithmetic mixing geometry of Collatz dynamics.

\textbf{Paper A} \cite{KamanoiA} established the renewal operator structure: the Syracuse map has a natural renewal decomposition, and the corresponding transfer operator has exact spectral properties.

\textbf{Paper B} \cite{KamanoiB} proved the exact TV mixing formula: for the Markov kernel $P_K$ on odd residues modulo $2^K$ induced by the 2-adic Haar measure, the support of $P_K^t(a,\cdot)$ forms an exact arithmetic progression, and
\[
  \TV(P_K^t(a,\cdot),\, \pi_K) = \max\!\left(0,\, 1 - 2^{S_t(a)-(K-1)}\right),
\]
yielding mixing time $T_{\mathrm{mix}}(K) = K-1$.

\textbf{Paper C} \cite{KamanoiC} introduced the one-bit high-bit kernel $P_{K,1}$ and observed the spectral jump $\delta_{K,0} = 0 \to \delta_{K,1} \approx 0.29$, formulating the Open Gap Problem: $\inf_K \delta_{K,1} > 0$.

The present paper asks: \emph{what do these population-level mixing results imply for individual Collatz orbits?}

\subsection{The fundamental gap}

Population-level mixing (Papers B and C) makes statements about distributions:
\begin{quote}
\emph{If $a$ is drawn uniformly at random from $\OK$, then $T^t(a) \bmod 2^K$ is close to uniform after $t = K-1$ steps.}
\end{quote}

The Collatz conjecture requires an individual-level statement:
\begin{quote}
\emph{For every fixed positive integer $n_0$, the orbit $n_0, T(n_0), T^2(n_0), \ldots$ eventually reaches $1$.}
\end{quote}

The gap between these levels is \emph{genuine}: a single deterministic orbit is not a random variable, and population mixing does not directly imply orbitwise convergence. Closing this gap is the central challenge of this paper.

\subsection{Main contributions}

\begin{enumerate}
\item \textbf{Effective high-bit parameter $m(t)$}: We introduce $m(t) = \lfloor\log_2 n_t\rfloor - K$ as the effective number of high bits active at time $t$, and show that for divergent orbits, $\limsup_t m(t) = \infty$ while $m(t) \to \infty$ drives the K-level dynamics toward the true Haar kernel---which has \emph{exponentially small} spectral gap (Section~\ref{sec:mt}).

\item \textbf{The self-mixing paradox}: Divergent orbits are in a sense \emph{poorly} mixing at each scale $K$, precisely because their large values push $m(t)$ into the regime where the Haar spectral gap is small. This reverses naive intuition (Section~\ref{sec:paradox}).

\item \textbf{Weyl equidistribution and its limits}: We connect to the known Weyl equidistribution result for Collatz orbits, and identify precisely why it does not suffice for the Collatz conjecture (Section~\ref{sec:weyl}).

\item \textbf{The Simultaneous Scale Coherence Problem}: We show the Collatz conjecture is equivalent to a multi-scale simultaneous condition, and formulate this as a precise open problem (Section~\ref{sec:ssc}).

\item \textbf{Conditional theorem}: If the Open Gap Problem of Paper C holds, then divergent orbits require an arithmetically impossible simultaneous coherence structure (Section~\ref{sec:conditional}).
\end{enumerate}

%% ============================================================
\section{The Effective High-Bit Parameter}
\label{sec:mt}
%% ============================================================

\begin{definition}[Effective high-bit parameter]
For an odd positive integer $n_t$ and a scale $K \geq 1$, the \emph{effective high-bit parameter at time $t$} is
\[
  m(t) = m(t; K) := \max\!\left(0,\; \lfloor \log_2 n_t \rfloor - K\right).
\]
\end{definition}

When $n_t \leq 2^K$, we have $m(t) = 0$ and the K-level dynamics is exactly deterministic: $n_{t+1} \bmod 2^K = F_K(n_t \bmod 2^K)$. When $n_t > 2^K$, the high bits of $n_t$ inject $m(t)$ bits of additional information into the K-level transition, and the K-level dynamics is governed by the kernel $P_{K,m(t)}$ of Paper C.

\begin{proposition}[High-bit inflow formula]
\label{prop:highbit}
For any odd $n_t$ with $m(t) \geq 1$:
\[
  n_{t+1} \bmod 2^K = F_K(n_t \bmod 2^K) + 3u_t \cdot 2^{K - v_t} \bmod 2^K,
\]
where $v_t = \vt(3n_t + 1)$ and $u_t = \lfloor n_t / 2^K \rfloor \bmod 2^{m(t)}$ encodes the active high bits.
\end{proposition}

\begin{proof}
Write $n_t = (n_t \bmod 2^K) + u_t \cdot 2^K$ with $0 \leq u_t < 2^{m(t)}$. By Lemma V-stability (Paper C, Lemma 2.3), $\vt(3n_t+1) = v_t = \vt(3(n_t \bmod 2^K)+1)$ when $v_t < K$. The formula follows by expanding $T(n_t) = (3n_t+1)/2^{v_t}$ and reducing modulo $2^K$.
\end{proof}

\subsection{Behavior for divergent orbits}

\begin{proposition}[Unbounded $m(t)$ for divergent orbits]
\label{prop:divergent-mt}
If $\{n_t\}$ is a divergent Collatz orbit ($n_t \to \infty$), then for every $K \geq 1$:
\[
  \limsup_{t \to \infty} m(t; K) = \infty.
\]
\end{proposition}

\begin{proof}
Since $n_t \to \infty$, for any $M$ there exist infinitely many $t$ with $n_t > 2^{K+M}$, so $m(t;K) > M$.
\end{proof}

The following table illustrates $m(t)$ for the orbit of $n_0 = 27$ at scale $K = 6$:

\begin{center}
\begin{tabular}{@{}rrrrc@{}}
\toprule
$t$ & $n_t$ & $n_t \bmod 2^6$ & $m(t)$ & $\delta_{K,m(t)}$ (approx.) \\
\midrule
 0 &   27 &  27 & 0 & $0$ \\
 4 &   71 &   7 & 0 & $0$ \\
 6 &  161 &  33 & 1 & $0.313$ \\
 9 &  137 &   9 & 1 & $0.313$ \\
14 &  263 &   7 & 2 & $0.428$ \\
16 &  593 &  17 & 3 & $0.571$ \\
\bottomrule
\end{tabular}
\end{center}

As $n_t$ grows, $m(t)$ increases and the effective spectral gap $\delta_{K,m(t)}$ increases---until the orbit falls back to small values.

%% ============================================================
\section{The Self-Mixing Paradox}
\label{sec:paradox}
%% ============================================================

\subsection{Large $m(t)$ does not imply fast mixing}

One might expect that a large $m(t)$---many high bits---implies fast K-level mixing. Paper C shows the opposite.

\begin{proposition}[Self-mixing paradox]
\label{prop:paradox}
As $m(t) \to \infty$, the effective kernel $P_{K,m(t)}$ approaches the true 2-adic Haar kernel $P_K^{\mathrm{Haar}}$, which has spectral gap $\delta_K^{\mathrm{Haar}} \approx 2^{-(K-1)} \to 0$ exponentially in $K$.
\end{proposition}

In other words: \emph{divergent orbits become progressively harder to mix spectrally at each scale $K$, because their large values push the K-level dynamics into the Haar regime with exponentially small spectral gap.}

This contrasts with the one-bit regime ($m=1$), where the spectral gap is $\delta_{K,1} \approx 0.29$ uniformly in $K$. The phase diagram (Paper C, Section 3) makes this precise:

\begin{center}
\begin{tabular}{@{}ccc@{}}
\toprule
$m(t)$ & Kernel regime & Spectral gap $\delta_{K,m}$ \\
\midrule
$0$ & Deterministic $F_K$ & $0$ \\
$1$ & One-bit flip & $\approx 0.29$ \\
$K$ & Arithmetic average & $\approx 0.74$ \\
$\infty$ (arith.) & Saturated average & $\approx 0.74$ \\
$\infty$ (true Haar) & 2-adic Haar & $\approx 2^{-(K-1)} \to 0$ \\
\bottomrule
\end{tabular}
\end{center}

The regime relevant for divergent orbits ($m(t) \to \infty$) is the bottom row: worst for spectral mixing.

\subsection{Interpretation}

Paper C showed: \emph{structured finite inflow beats full Haar randomization}. The present paper refines this: divergent orbits are in the Haar regime, which is the least efficient for spectral decorrelation. This means divergent orbits \emph{self-inhibit} their own mixing---a structural constraint that will be exploited in the conditional theorem (Section~\ref{sec:conditional}).

%% ============================================================
\section{Weyl Equidistribution and Its Limits}
\label{sec:weyl}
%% ============================================================

\subsection{What equidistribution gives}

A classical result (Terras \cite{Terras1976}, Lagarias \cite{Lagarias1985}) states that Collatz orbits are \emph{equidistributed} modulo $2^K$:

\begin{theorem}[Weyl equidistribution, after Terras--Lagarias]
\label{thm:weyl}
For any positive odd integer $n_0$ and any odd residue $r \in \OK$:
\[
  \lim_{T \to \infty} \frac{1}{T} \#\{t \leq T : n_t \equiv r \pmod{2^K}\} = \frac{1}{n_K},
\]
assuming the orbit $\{n_t\}$ is infinite (i.e., assuming $n_0$ is divergent or the orbit is long before reaching $1$).
\end{theorem}

In particular, for $r = 1$: every infinite orbit visits $n_t \equiv 1 \pmod{2^K}$ with limiting frequency $1/n_K$.

\subsection{Why this does not suffice}

The condition $n_t \equiv 1 \pmod{2^K}$ means $n_t \in \{1, 1+2^K, 1+2\cdot 2^K, \ldots\}$. For a divergent orbit, only values $n_t = 1 + j \cdot 2^K$ with $j \geq 1$ are visited (not $n_t = 1$ itself). The equidistribution theorem does not distinguish between $n_t = 1$ and $n_t = 1 + j \cdot 2^K$ for $j \geq 1$.

\begin{remark}[The equidistribution gap]
Theorem~\ref{thm:weyl} gives frequency $1/n_K$ of visiting $n_t \equiv 1 \pmod{2^K}$, but does \emph{not} imply that $n_t = 1$ is ever reached. A divergent orbit can satisfy the equidistribution theorem at every scale $K$ while never actually reaching $1$.
\end{remark}

This gap is not a deficiency of the argument; it is intrinsic to frequency-based methods applied to the integer $1$.

%% ============================================================
\section{The Simultaneous Scale Coherence Problem}
\label{sec:ssc}
%% ============================================================

\subsection{2-adic characterization of convergence}

The key observation is algebraic: the positive integer $1$ is the \emph{unique} positive integer satisfying
\[
  n \equiv 1 \pmod{2^K} \quad\text{for all } K \geq 1.
\]
This is because the 2-adic integers $\mathbb{Z}_2$ embed $\mathbb{Z}_{>0}$ via the map $n \mapsto (n \bmod 2^K)_{K \geq 1}$, and the 2-adic representation of $1$ is $(1, 1, 1, \ldots)$. Any other positive integer $m \geq 3$ fails $m \equiv 1 \pmod{2^K}$ for $K = \lfloor \log_2(m-1) \rfloor + 1$.

\begin{proposition}[2-adic criterion for convergence]
\label{prop:2adic}
A Collatz orbit $\{n_t\}$ converges to $1$ if and only if:
\[
  \exists t_0 : \forall K \geq 1,\quad n_{t_0} \equiv 1 \pmod{2^K}.
\]
Equivalently: $\exists t_0 : n_{t_0} = 1$.
\end{proposition}

\subsection{The simultaneous scale condition}

\begin{definition}[Scale-coherence at $1$]
An orbit $\{n_t\}$ is \emph{$(K,R)$-coherent at $1$} if there exists $t_0 \leq R$ such that $n_{t_0} \equiv 1 \pmod{2^K}$.

It is \emph{$(\infty,R)$-coherent at $1$} if there exists $t_0 \leq R$ such that $n_{t_0} \equiv 1 \pmod{2^K}$ for \textbf{all} $K \geq 1$ simultaneously.
\end{definition}

By Proposition~\ref{prop:2adic}, $(\infty,R)$-coherence at $1$ is equivalent to $n_{t_0} = 1$ for some $t_0 \leq R$.

\begin{openproblem}[Simultaneous Scale Coherence Problem]
\label{prob:ssc}
Let $\{n_t\}$ be any Collatz orbit. Prove:
\[
  \forall K \geq 1,\; \exists t_0 : n_{t_0} \equiv 1 \pmod{2^K} \quad\Longrightarrow\quad \exists t_0 : n_{t_0} = 1.
\]
\end{openproblem}

\begin{remark}
The left-hand side of Open Problem~\ref{prob:ssc} is a consequence of Theorem~\ref{thm:weyl}: by equidistribution, every infinite orbit visits $n_t \equiv 1 \pmod{2^K}$ for each fixed $K$. The challenge is the \emph{simultaneity}: finding a single $t_0$ that works for all $K$ at once.
\end{remark}

\begin{remark}[Equivalence to Collatz]
Open Problem~\ref{prob:ssc} is equivalent to the Collatz conjecture. The implication ``$\exists t_0: n_{t_0} = 1$ $\Rightarrow$ each-$K$ condition'' is trivial. The converse is what needs proof.
\end{remark}

The following table illustrates how the constraint tightens with $K$:

\begin{center}
\begin{tabular}{@{}rrr@{}}
\toprule
$K$ & Mod $2^K$ constraint & Minimum $n > 1$ satisfying it \\
\midrule
 1 &  $n \equiv 1 \pmod{2}$   & $3$ \\
 4 &  $n \equiv 1 \pmod{16}$  & $17$ \\
 8 &  $n \equiv 1 \pmod{256}$ & $257$ \\
12 &  $n \equiv 1 \pmod{4096}$ & $4097$ \\
$K$ &  $n \equiv 1 \pmod{2^K}$ & $1 + 2^K$ \\
\bottomrule
\end{tabular}
\end{center}

A divergent orbit satisfying the constraint at every $K$ simultaneously would need $n_{t_0} > 2^K$ for all $K$---impossible for any finite $n_{t_0}$.

%% ============================================================
\section{Conditional Theorem}
\label{sec:conditional}
%% ============================================================

We now state the main conditional result, connecting the Open Gap Problem (Paper C) to the Simultaneous Scale Coherence Problem.

\subsection{Setup}

Let $\{n_t\}$ be a Collatz orbit. At time $t$ and scale $K$, the effective kernel is $P_{K,m(t)}$ with spectral gap $\delta_{K,m(t)}$. Define the \emph{cumulative spectral mixing} up to time $T$ at scale $K$ as:
\[
  \Delta_K(T) := \sum_{t=1}^{T} \delta_{K,m(t)}.
\]

\begin{theorem}[Conditional Scale-Coherence Rigidity]
\label{thm:conditional}
Assume the Open Gap Problem of Paper C holds: $\inf_K \delta_{K,1} \geq c > 0$.

Let $\{n_t\}$ be a Collatz orbit with $\liminf_{t \to \infty} n_t > 2^K$. Then:
\[
  \Delta_K(T) \geq c \cdot \#\{t \leq T : n_t > 2^K\} \to \infty \quad\text{as } T \to \infty.
\]

In particular, the K-level spectral mixing accumulated by the orbit diverges, implying that the K-level distribution of $n_t \bmod 2^K$ becomes asymptotically indistinguishable from the stationary distribution $\pi_{K,m}$.
\end{theorem}

\begin{proof}
For each $t$ with $n_t > 2^K$, we have $m(t;K) \geq 1$, so $\delta_{K,m(t)} \geq \delta_{K,1} \geq c$. Summing over all such $t$ gives $\Delta_K(T) \geq c \cdot \#\{t \leq T : n_t > 2^K\}$. If $\liminf n_t > 2^K$, then almost all $t$ satisfy $n_t > 2^K$, so the count grows as $T$.
\end{proof}

\subsection{Interpretation and the remaining gap}

Theorem~\ref{thm:conditional} establishes that divergent orbits accumulate diverging spectral mixing at each fixed scale $K$. This means the K-level time-averaged behavior approaches stationarity.

However---and this is the honest statement of the remaining gap---accumulated spectral mixing in the \emph{population} (or time-averaged) sense does not directly force the \emph{individual} orbit to hit $n_t = 1$. The theorem shows increasing decorrelation in the time average, but the orbit remains a single deterministic trajectory.

\begin{remark}[The remaining gap]
Closing the gap between Theorem~\ref{thm:conditional} and a proof of the Collatz conjecture requires showing that diverging cumulative spectral mixing at \emph{all} scales simultaneously forces a single time $t_0$ where $n_{t_0} = 1$. This is precisely the Simultaneous Scale Coherence Problem (Open Problem~\ref{prob:ssc}).
\end{remark}

The contribution of this paper is to make the gap explicit and to show that it is not a vague philosophical distinction but a precise mathematical problem: the transition from ``each-$K$ equidistribution'' to ``all-$K$ simultaneously''.

%% ============================================================
\section{The Structure of Papers B, C, D}
\label{sec:structure}
%% ============================================================

We now articulate the logical structure of the series, showing how Papers B, C, D form a coherent program.

\begin{center}
\renewcommand{\arraystretch}{1.4}
\begin{tabular}{@{}p{1.8cm}p{4.5cm}p{4.5cm}@{}}
\toprule
Paper & Core result & Status \\
\midrule
\textbf{B} & Exact TV mixing: $\TV(P_K^t, \pi_K) = \max(0, 1-2^{S_t-(K-1)})$ & Proved \\
\textbf{C} & One-bit spectral jump: $\delta_{K,0}=0 \to \delta_{K,1}\approx 0.29$ & Numerical; Open Gap Problem \\
\textbf{D} & Simultaneous scale coherence: divergence requires impossible memory & Open; conditional on C \\
\bottomrule
\end{tabular}
\end{center}

The logical dependencies are:

\[
  \underbrace{\text{Paper B (ETV)}}_{\text{proved}} \;\longrightarrow\; \underbrace{\text{Theorem D (conditional)}}_{\text{conditional on C}} \;\longleftarrow\; \underbrace{\text{Paper C (Open Gap)}}_{\text{open}}.
\]

If Paper C's Open Gap Problem is resolved, Theorem~\ref{thm:conditional} gives diverging spectral mixing for all orbits with $\liminf n_t > 2^K$. The remaining bridge---Simultaneous Scale Coherence (Open Problem~\ref{prob:ssc})---is equivalent to the Collatz conjecture itself.

%% ============================================================
\section{Summary of Results}
%% ============================================================

\begin{center}
\begin{tabular}{@{}lll@{}}
\toprule
Result & Status & Section \\
\midrule
Effective high-bit parameter $m(t)$ & Definition & \S\ref{sec:mt} \\
Proposition: $\limsup m(t) = \infty$ for divergent orbits & Proved & \S\ref{sec:mt} \\
Self-mixing paradox (large $m \Rightarrow$ small gap) & Proved & \S\ref{sec:paradox} \\
Weyl equidistribution and its limits & Known/Explained & \S\ref{sec:weyl} \\
2-adic criterion for convergence & Proved & \S\ref{sec:ssc} \\
Simultaneous Scale Coherence Problem & \textbf{Open} (= Collatz) & \S\ref{sec:ssc} \\
Conditional Theorem (Open Gap $\Rightarrow$ diverging mixing) & Conditional & \S\ref{sec:conditional} \\
The remaining gap (population $\not\Rightarrow$ individual) & Identified & \S\ref{sec:conditional} \\
\bottomrule
\end{tabular}
\end{center}

\begin{thebibliography}{99}

\bibitem{KamanoiA}
H. Kamanoi,
``Exact Operator Structure and Finite Information Memory in Accelerated Collatz Dynamics'' (Paper A), 2025.

\bibitem{KamanoiB}
H. Kamanoi,
``Exact Finite-Level Mixing Geometry in Collatz Dynamics'' (Paper B), 2026.

\bibitem{KamanoiC}
H. Kamanoi,
``Spectral Amplification by One-Bit Inflow in Collatz Dynamics'' (Paper C), 2026.

\bibitem{Lagarias1985}
J.~C. Lagarias, ``The $3x+1$ problem and its generalizations,''
\textit{Amer.\ Math.\ Monthly} \textbf{92} (1985), 3--23.

\bibitem{Lagarias2010}
J.~C. Lagarias (ed.), \textit{The Ultimate Challenge: The $3x+1$ Problem}, AMS, 2010.

\bibitem{Terras1976}
R. Terras, ``A stopping time problem on the positive integers,''
\textit{Acta Arith.} \textbf{30} (1976), 241--252.

\end{thebibliography}

\end{document}
